\IfFileExists{stacks-project.cls}{%
\documentclass{stacks-project}
}{%
\documentclass{amsart}
}
\usepackage{amssymb}

% For dealing with references we use the comment environment
\usepackage{verbatim}
\usepackage{fancyvrb}
\newenvironment{reference}{\comment}{\endcomment}
%\newenvironment{reference}{}{}
\newenvironment{slogan}{\comment}{\endcomment}
\newenvironment{history}{\comment}{\endcomment}

% For source figures retained as assets
\usepackage{graphicx}

% For commutative diagrams we use Xy-pic
\usepackage[all]{xy}

% We use 2cell for 2-commutative diagrams.
\xyoption{2cell}
\UseAllTwocells

% We use multicol for the list of chapters between chapters
\usepackage{multicol}

% This is generally recommended for better output
\usepackage{lmodern}
\usepackage[T1]{fontenc}
\usepackage{textalpha}

% For cross-file-references
% We prefer the snapshot tmp/xref/#2.aux,
% which latexmkrc replaces atomically, so
% that documents compiled simultaneously
% never read a half-written .aux file.
% The last argument keeps links to #2.pdf.
\usepackage{xr-hyper}
\newcommand{\maybeexternaldocument}[2]{%
  \IfFileExists{#2.aux}{%
    \externaldocument[#1]{#2}%
  }{%
    \PackageWarning{stack}{Missing #2.aux; external references with prefix #1 unavailable}%
  }%
}
\let\externaldocumentorig\externaldocument
\renewcommand{\externaldocument}[2][]{%
  \IfFileExists{tmp/xref/#2.aux}{%
    \externaldocumentorig[#1]{tmp/xref/#2}[#2.pdf]%
  }{%
  \IfFileExists{#2.aux}{%
    \externaldocumentorig[#1]{#2}%
  }{%
    \PackageWarning{stack}{Missing #2.aux; external references with prefix #1 unavailable}%
  }%
  }%
}

% Package for hypertext links:
\usepackage{hyperref}

% For any local file, say "hello.tex" you want to link to please
% use \externaldocument[hello-]{hello}
\externaldocument[introduction-]{introduction}
\externaldocument[conventions-]{conventions}
\externaldocument[sets-]{sets}
\externaldocument[categories-]{categories}
\externaldocument[topology-]{topology}
\externaldocument[sheaves-]{sheaves}
\externaldocument[sites-]{sites}
\externaldocument[stacks-]{stacks}
\externaldocument[fields-]{fields}
\externaldocument[algebra-]{algebra}
\externaldocument[brauer-]{brauer}
\externaldocument[homology-]{homology}
\externaldocument[derived-]{derived}
\externaldocument[simplicial-]{simplicial}
\externaldocument[more-algebra-]{more-algebra}
\externaldocument[smoothing-]{smoothing}
\externaldocument[modules-]{modules}
\externaldocument[sites-modules-]{sites-modules}
\externaldocument[injectives-]{injectives}
\externaldocument[cohomology-]{cohomology}
\externaldocument[sites-cohomology-]{sites-cohomology}
\externaldocument[dga-]{dga}
\externaldocument[dpa-]{dpa}
\externaldocument[sdga-]{sdga}
\externaldocument[hypercovering-]{hypercovering}
\externaldocument[schemes-]{schemes}
\externaldocument[constructions-]{constructions}
\externaldocument[properties-]{properties}
\externaldocument[morphisms-]{morphisms}
\externaldocument[coherent-]{coherent}
\externaldocument[divisors-]{divisors}
\externaldocument[limits-]{limits}
\externaldocument[varieties-]{varieties}
\externaldocument[topologies-]{topologies}
\externaldocument[descent-]{descent}
\externaldocument[perfect-]{perfect}
\externaldocument[more-morphisms-]{more-morphisms}
\externaldocument[flat-]{flat}
\externaldocument[groupoids-]{groupoids}
\externaldocument[more-groupoids-]{more-groupoids}
\externaldocument[etale-]{etale}
\externaldocument[chow-]{chow}
\externaldocument[intersection-]{intersection}
\externaldocument[pic-]{pic}
\externaldocument[weil-]{weil}
\externaldocument[adequate-]{adequate}
\externaldocument[dualizing-]{dualizing}
\externaldocument[duality-]{duality}
\externaldocument[discriminant-]{discriminant}
\externaldocument[derham-]{derham}
\externaldocument[local-cohomology-]{local-cohomology}
\externaldocument[algebraization-]{algebraization}
\externaldocument[curves-]{curves}
\externaldocument[resolve-]{resolve}
\externaldocument[models-]{models}
\externaldocument[functors-]{functors}
\externaldocument[equiv-]{equiv}
\externaldocument[pione-]{pione}
\externaldocument[etale-cohomology-]{etale-cohomology}
\externaldocument[proetale-]{proetale}
\externaldocument[relative-cycles-]{relative-cycles}
\externaldocument[more-etale-]{more-etale}
\externaldocument[trace-]{trace}
\externaldocument[crystalline-]{crystalline}
\externaldocument[spaces-]{spaces}
\externaldocument[spaces-properties-]{spaces-properties}
\externaldocument[spaces-morphisms-]{spaces-morphisms}
\externaldocument[decent-spaces-]{decent-spaces}
\externaldocument[spaces-cohomology-]{spaces-cohomology}
\externaldocument[spaces-limits-]{spaces-limits}
\externaldocument[spaces-divisors-]{spaces-divisors}
\externaldocument[spaces-over-fields-]{spaces-over-fields}
\externaldocument[spaces-topologies-]{spaces-topologies}
\externaldocument[spaces-descent-]{spaces-descent}
\externaldocument[spaces-perfect-]{spaces-perfect}
\externaldocument[spaces-more-morphisms-]{spaces-more-morphisms}
\externaldocument[spaces-flat-]{spaces-flat}
\externaldocument[spaces-groupoids-]{spaces-groupoids}
\externaldocument[spaces-more-groupoids-]{spaces-more-groupoids}
\externaldocument[bootstrap-]{bootstrap}
\externaldocument[spaces-pushouts-]{spaces-pushouts}
\externaldocument[spaces-chow-]{spaces-chow}
\externaldocument[groupoids-quotients-]{groupoids-quotients}
\externaldocument[spaces-more-cohomology-]{spaces-more-cohomology}
\externaldocument[spaces-simplicial-]{spaces-simplicial}
\externaldocument[spaces-duality-]{spaces-duality}
\externaldocument[formal-spaces-]{formal-spaces}
\externaldocument[restricted-]{restricted}
\externaldocument[spaces-resolve-]{spaces-resolve}
\externaldocument[formal-defos-]{formal-defos}
\externaldocument[defos-]{defos}
\externaldocument[cotangent-]{cotangent}
\externaldocument[examples-defos-]{examples-defos}
\externaldocument[algebraic-]{algebraic}
\externaldocument[examples-stacks-]{examples-stacks}
\externaldocument[stacks-sheaves-]{stacks-sheaves}
\externaldocument[criteria-]{criteria}
\externaldocument[artin-]{artin}
\externaldocument[quot-]{quot}
\externaldocument[stacks-properties-]{stacks-properties}
\externaldocument[stacks-morphisms-]{stacks-morphisms}
\externaldocument[stacks-limits-]{stacks-limits}
\externaldocument[stacks-cohomology-]{stacks-cohomology}
\externaldocument[stacks-perfect-]{stacks-perfect}
\externaldocument[stacks-introduction-]{stacks-introduction}
\externaldocument[stacks-more-morphisms-]{stacks-more-morphisms}
\externaldocument[stacks-geometry-]{stacks-geometry}
\externaldocument[moduli-]{moduli}
\externaldocument[moduli-curves-]{moduli-curves}
\externaldocument[examples-]{examples}
\externaldocument[exercises-]{exercises}
\externaldocument[guide-]{guide}
\externaldocument[desirables-]{desirables}
\externaldocument[coding-]{coding}
\externaldocument[obsolete-]{obsolete}
\externaldocument[fdl-]{fdl}
\externaldocument[index-]{index}

% Heap Project documents
\externaldocument[k3-]{k3}
\externaldocument[torelli-]{torelli}
\externaldocument[linear-algebra-]{linear-algebra}
\externaldocument[groups-]{groups}
\externaldocument[complex-analysis-]{complex-analysis}
\externaldocument[functional-analysis-]{functional-analysis}
\externaldocument[categories-]{categories}
\externaldocument[commutative-algebra-]{commutative-algebra}
\externaldocument[ringed-spaces-]{ringed-spaces}
\externaldocument[homological-algebra-]{homological-algebra}
\externaldocument[lie-algebras-]{lie-algebras}
\externaldocument[representation-theory-]{representation-theory}
\externaldocument[smooth-manifolds-]{smooth-manifolds}
\externaldocument[differential-topology-]{differential-topology}
\externaldocument[algebraic-topology-]{algebraic-topology}
\externaldocument[differential-geometry-]{differential-geometry}
\externaldocument[algebraic-geometry-]{algebraic-geometry}
\externaldocument[symplectic-geometry-]{symplectic-geometry}
\externaldocument[complex-geometry-]{complex-geometry}
\externaldocument[hodge-theory-]{hodge-theory}
\externaldocument[%
variations-of-hodge-structures-]{%
variations-of-hodge-structures}
\externaldocument[%
homological-mirror-symmetry-]{%
homological-mirror-symmetry}
\maybeexternaldocument{deformations-}{deformations}
\externaldocument[vertex-algebras-]{vertex-algebras}
\externaldocument[springer-]{springer}
\externaldocument[infty-categories-]{infty-categories}
\externaldocument[seminars-differential-geometry-]{%
seminars-differential-geometry}
\externaldocument[seminars-complex-geometry-]{%
seminars-complex-geometry}
\externaldocument[seminars-algebraic-geometry-]{%
seminars-algebraic-geometry}
\externaldocument[seminars-others-]{%
seminars-others}
\externaldocument[%
geometric-prequantization-]{%
geometric-prequantization}
\externaldocument[surfaces-]{surfaces}
\externaldocument[birational-maps-commutative-algebra-]{birational-maps-commutative-algebra}
\externaldocument[cremona-transformations-]{cremona-transformations}
\externaldocument[derived-categories-of-sheaves-]{derived-categories-of-sheaves}
\externaldocument[geometric-invariant-theory-]{geometric-invariant-theory}
\externaldocument[geometric-stability-conditions-and-group-actions-]{geometric-stability-conditions-and-group-actions}
\externaldocument[moduli-spaces-of-sheaves-]{moduli-spaces-of-sheaves}
\externaldocument[bridgeland-stability-general-theory-]{bridgeland-stability-general-theory}
\externaldocument[hilbert-schemes-]{hilbert-schemes}
\externaldocument[noncommutative-abelian-surfaces-and-kummer-type-hyperkahler-varieties-]{noncommutative-abelian-surfaces-and-kummer-type-hyperkahler-varieties}
\externaldocument[mumford-tate-groups-in-hodge-theory-]{mumford-tate-groups-in-hodge-theory}
\externaldocument[hodge-birational-atoms-2026-]{hodge-birational-atoms-2026}
\externaldocument[xxii-egd-2026-]{%
xxii-egd-2026}
\externaldocument[pde-]{pde}
\externaldocument[probability-]{probability}

\externaldocument[linguistics-]{linguistics}
\externaldocument[genealogia-familiar-]{%
genealogia-familiar}

\externaldocument[%
16th-alga-meeting-2026-]{%
16th-alga-meeting-2026}

\externaldocument[%
iv-jornadas-lefschetz-2026-]{%
iv-jornadas-lefschetz-2026}

\externaldocument[reading-the-masters-]{%
reading-the-masters}

\externaldocument[real-analysis-]{%
real-analysis}

\externaldocument[optimal-transport-]{%
optimal-transport}

\externaldocument[history-]{history}

\externaldocument[seiberg-witten-]{seiberg-witten}

\externaldocument[history-philosophy-biology-]{history-philosophy-biology}

\externaldocument[chemistry-]{chemistry}

\externaldocument[molecular-cell-biology-]{molecular-cell-biology}

\externaldocument[microbiology-]{microbiology}

\externaldocument[protistology-]{protistology}

\externaldocument[mycology-]{mycology}

\externaldocument[botany-]{botany}

\externaldocument[zoology-]{zoology}

\externaldocument[animal-histology-]{animal-histology}

\externaldocument[animal-physiology-]{animal-physiology}

\externaldocument[systematics-biogeography-]{systematics-biogeography}

\externaldocument[ecology-natural-resources-]{ecology-natural-resources}

\externaldocument[earth-sciences-]{earth-sciences}

\externaldocument[biostatistics-]{biostatistics}

\externaldocument[mathematical-systems-biology-]{mathematical-systems-biology}

\externaldocument[cimpa-el-salvador-2024-]{cimpa-el-salvador-2024}

% Theorem environments.
%
\theoremstyle{plain}
\newtheorem{theorem}[subsection]{Theorem}
\newtheorem{proposition}[subsection]{Proposition}
\newtheorem{lemma}[subsection]{Lemma}
\newtheorem{corollary}[subsection]{Corollary}

\theoremstyle{definition}
\newtheorem{definition}[subsection]{Definition}
\newtheorem{example}[subsection]{Example}
\newtheorem{exercise}[subsection]{Exercise}
\newtheorem{situation}[subsection]{Situation}

\theoremstyle{remark}
\newtheorem{remark}[subsection]{Remark}
\newtheorem{remarks}[subsection]{Remarks}

\numberwithin{equation}{subsection}

% Macros
%
\def\lim{\mathop{\mathrm{lim}}\nolimits}
\def\colim{\mathop{\mathrm{colim}}\nolimits}
\def\Spec{\mathop{\mathrm{Spec}}}
\def\Hom{\mathop{\mathrm{Hom}}\nolimits}
\def\Ext{\mathop{\mathrm{Ext}}\nolimits}
\def\SheafHom{\mathop{\mathcal{H}\!\mathit{om}}\nolimits}
\def\SheafExt{\mathop{\mathcal{E}\!\mathit{xt}}\nolimits}
\def\Sch{\mathit{Sch}}
\def\Mor{\mathop{\mathrm{Mor}}\nolimits}
\def\Ob{\mathop{\mathrm{Ob}}\nolimits}
\def\Sh{\mathop{\mathit{Sh}}\nolimits}
\def\NL{\mathop{N\!L}\nolimits}
\def\CH{\mathop{\mathrm{CH}}\nolimits}
\def\proetale{{pro\text{-}\acute{e}tale}}
\def\etale{{\acute{e}tale}}
\def\QCoh{\mathit{QCoh}}
\def\Ker{\mathop{\mathrm{Ker}}}
\def\Im{\mathop{\mathrm{Im}}}
\def\Coker{\mathop{\mathrm{Coker}}}
\def\Coim{\mathop{\mathrm{Coim}}}

% Boxtimes
%
\DeclareMathSymbol{\boxtimes}{\mathbin}{AMSa}{"02}

%
% Macros for moduli stacks/spaces
%
\def\QCohstack{\mathcal{QC}\!\mathit{oh}}
\def\Cohstack{\mathcal{C}\!\mathit{oh}}
\def\Spacesstack{\mathcal{S}\!\mathit{paces}}
\def\Quotfunctor{\mathrm{Quot}}
\def\Hilbfunctor{\mathrm{Hilb}}
\def\Curvesstack{\mathcal{C}\!\mathit{urves}}
\def\Polarizedstack{\mathcal{P}\!\mathit{olarized}}
\def\Complexesstack{\mathcal{C}\!\mathit{omplexes}}
% \Pic is the operator that assigns to X its picard group, usage \Pic(X)
% \Picardstack_{X/B} denotes the Picard stack of X over B
% \Picardfunctor_{X/B} denotes the Picard functor of X over B
\def\Pic{\mathop{\mathrm{Pic}}\nolimits}
\def\Picardstack{\mathcal{P}\!\mathit{ic}}
\def\Picardfunctor{\mathrm{Pic}}
\def\Deformationcategory{\mathcal{D}\!\mathit{ef}}


\begin{document}

\title{Biostatistics}
\maketitle

\phantomsection
\label{section-phantom}

\tableofcontents

\section{What this chapter covers}
\label{section-overview}

Classical statistics for biologists:
probability and the binomial and Poisson
distributions, confidence intervals,
hypothesis tests, comparison of two or
more groups, simple linear regression,
the chi-square test of independence, and
sampling designs.

\textbf{Sources of these notes.}
The notes rest on the course
\emph{Bioestad\'istica}, third semester,
Fall 2014 (semester 2015-1), Facultad de
Ciencias, UNAM. What survives is
Daniel's take-home probability exam,
exercise sheets on confidence intervals,
tests and ANOVA with his solutions, JMP
outputs of his analyses, Excel workbooks
on regression and the chi-square test,
the teacher's slides on sampling, and an
exam answer on a salamander survey. The
course used JMP and Excel and cites no
textbook. Definitions and formulas were
checked against Rosner
\cite{rosner2011biostatistics} and Triola
\cite{triola2012elementary}, and all
numbers below were recomputed.

\section{Describing a sample}
\label{section-descriptive}

Every analysis in the course started
with a statistical summary of each
sample: $n$, mean, standard deviation,
standard error, variance, skewness,
kurtosis, coefficient of variation,
quartiles, minimum and maximum, then a
histogram and a box plot. For a sample
$x_1, \dots, x_n$:
$$
\bar{x} = \frac{1}{n}\sum_{i=1}^n x_i,
\qquad
s^2 = \frac{1}{n-1}
\sum_{i=1}^n (x_i - \bar{x})^2,
$$
$$
\mathrm{SE} = \frac{s}{\sqrt{n}},
\qquad
\mathrm{CV} = 100\,\frac{s}{\bar{x}}\ \%.
$$
The standard deviation describes the
spread of the data; the standard error
describes the uncertainty of the mean
\cite[Section 2.6]{rosner2011biostatistics},
\cite[Section
6.5]{rosner2011biostatistics}.

\section{Probability}
\label{section-probability}

For events $A$, $B$ with $P(B) > 0$ the
conditional probability is
$$
P(A \mid B) =
\frac{P(A \cap B)}{P(B)}.
$$
Two events are mutually exclusive if
$P(A \cap B) = 0$, and independent if
$$
P(A \cap B) = P(A)\,P(B),
$$
equivalently $P(A \mid B) = P(A)$.
Mutually exclusive events with positive
probabilities are never independent. If
$B_1, \dots, B_k$ are mutually exclusive
and exhaustive, the law of total
probability and Bayes' rule read
\begin{equation}
\label{equation-bayes}
P(B_j \mid A) =
\frac{P(A \mid B_j)\,P(B_j)}
{\sum_{i=1}^k P(A \mid B_i)\,P(B_i)}.
\end{equation}
Here $P(B_j)$ is the prior (\emph{a
priori}) and $P(B_j \mid A)$ the
posterior (\emph{a posteriori})
probability
\cite[Sections
3.4--3.7]{rosner2011biostatistics}.

\begin{example}
\label{example-bayes-flowers}
In the take-home exam, two plant
populations of $100$ individuals each
carry a flower-colour gene: white,
yellow and orange flowers number
$60, 24, 16$ in the first population and
$32, 48, 20$ in the second. A flower is
drawn from a population chosen with
probability $1/2$ and turns out yellow.
By \eqref{equation-bayes}
$$
P(P_2 \mid Y) =
\frac{\frac{1}{2}(0.48)}
{\frac{1}{2}(0.48) + \frac{1}{2}(0.24)}
= \frac{2}{3}.
$$
\end{example}

\section{Binomial and Poisson distributions}
\label{section-binomial-poisson}

A random variable $X$ counting successes
in $n$ independent trials, each with
success probability $p$, is binomial:
$$
P(X = x) = \binom{n}{x}
p^x (1-p)^{n-x},
\qquad x = 0, \dots, n,
$$
with mean $np$ and variance $np(1-p)$. A
count of rare events in a fixed interval
of time or space, with mean $\lambda$, is
Poisson:
$$
P(X = x) = \frac{e^{-\lambda}
\lambda^x}{x!},
\qquad x = 0, 1, 2, \dots,
$$
with mean and variance both equal to
$\lambda$. The binomial needs the number
of trials; the Poisson needs only the
mean rate
\cite[Sections
4.8--4.12]{rosner2011biostatistics}.

\begin{example}
\label{example-binomial-cables}
In an exam problem each of $15$
undersea cables is bitten by sharks with
probability $0.35$, so $X$ is binomial
with $n = 15$ and $p = 0.35$. Summing the
probability function:
\begin{center}
\begin{tabular}{lc}
Event & Probability \\
\hline
$X = 0$ & $0.0016$ \\
$X = 4$ & $0.1792$ \\
$X \leq 4$ & $0.3519$ \\
$X \geq 6$ & $0.4357$ \\
$5 \leq X \leq 9$ & $0.6356$ \\
\end{tabular}
\end{center}
In another problem a hitchhiking robot is
picked up with probability $0.2$ on each
of $27$ attempts: the number of lifts has
mean $27(0.2) = 5.4$, variance
$27(0.2)(0.8) = 4.32$ and standard
deviation $2.08$.
\end{example}

\begin{example}
\label{example-poisson-larvae}
Beetle larvae per soil sample follow a
Poisson law with $\lambda = 4.2$. Then
$P(X = 4) = 0.194$,
$P(X < 3) = 0.210$ and
$P(X > 4) = 0.410$.
\end{example}

\section{Confidence intervals}
\label{section-confidence-intervals}

Write $z_q$, $t_{d,q}$ and $\chi^2_{d,q}$
for the $q$-quantiles of the standard
normal, the $t$ distribution and the
chi-square distribution with $d$ degrees
of freedom. A $100(1-\alpha)\%$
confidence interval is:
\begin{itemize}
\item for a mean, with known $\sigma$ or
large $n$:
$$
\bar{x} \pm z_{1-\alpha/2}
\frac{\sigma}{\sqrt{n}};
$$
\item for a mean of a normal population
with unknown $\sigma$:
$$
\bar{x} \pm t_{n-1,1-\alpha/2}
\frac{s}{\sqrt{n}};
$$
\item for a variance of a normal
population:
$$
\left[
\frac{(n-1)s^2}{\chi^2_{n-1,1-\alpha/2}},
\ \frac{(n-1)s^2}{\chi^2_{n-1,\alpha/2}}
\right];
$$
\item for a proportion, with
$n\hat{p}(1-\hat{p}) \geq 5$:
$$
\hat{p} \pm z_{1-\alpha/2}
\sqrt{\frac{\hat{p}(1-\hat{p})}{n}}.
$$
\end{itemize}
See \cite[Sections
6.5--6.8]{rosner2011biostatistics}. The
interval is random; the population
parameter is fixed. In repeated sampling
$100(1-\alpha)\%$ of such intervals
contain it.

\begin{example}
\label{example-confidence-intervals}
From the exercise sheet:
\begin{itemize}
\item $\bar{x} = 817.75$, $s = 73.62$,
$n = 30$, $99\%$: with
$z_{0.995} = 2.576$ and
$\mathrm{SE} = 13.44$ the interval is
$783.13 \leq \mu \leq 852.37$.
\item $s = 29.3$, $n = 12$, $99\%$: with
$\chi^2_{11,0.995} = 26.76$ and
$\chi^2_{11,0.005} = 2.603$ the interval
is $352.9 \leq \sigma^2 \leq 3627.6$,
that is $18.8 \leq \sigma \leq 60.2$.
\item $\hat{p} = 0.51$, $n = 60$, $95\%$:
$0.384 \leq p \leq 0.636$.
\end{itemize}
\end{example}

\section{Hypothesis tests}
\label{section-tests}

A test compares a null hypothesis $H_0$
with an alternative $H_1$. A test
statistic is computed from the data; the
$p$-value is the probability, under
$H_0$, of a value at least as extreme as
the one observed. One rejects $H_0$ at
level $\alpha$ when $p < \alpha$,
equivalently when the statistic falls
in the rejection region. A type I error
rejects a true $H_0$ (probability
$\alpha$); a type II error accepts a
false one. A one-sided alternative puts
the whole rejection region in one tail;
it is legitimate only when the direction
is fixed before seeing the data
\cite[Sections
7.2--7.4]{rosner2011biostatistics}.

\textbf{Homoscedasticity.}
The pooled $t$ test and ANOVA assume equal
variances. For two normal samples the $F$
test uses
$$
F = \frac{s_1^2}{s_2^2}
\sim F_{n_1-1,\,n_2-1}
$$
under $H_0\colon \sigma_1^2 = \sigma_2^2$
\cite[Section
8.6]{rosner2011biostatistics}.
It is sensitive to non-normality. The
Levene test (and its Brown-Forsythe
version, with deviations from the median)
replaces each value by its absolute
deviation from the group mean or median
and compares these deviations across
groups \cite[Section
9-5]{triola2012elementary}. JMP also
reports Bartlett's and O'Brien's tests;
Daniel's analyses checked all four
before each comparison of means.

\section{Comparing two groups}
\label{section-two-samples}

For independent normal samples with
equal variances, the pooled $t$
statistic is
$$
t = \frac{\bar{x}_1 - \bar{x}_2}
{s_p\sqrt{\frac{1}{n_1}+\frac{1}{n_2}}},
\qquad
s_p^2 = \frac{(n_1-1)s_1^2
+ (n_2-1)s_2^2}{n_1+n_2-2},
$$
with $n_1 + n_2 - 2$ degrees of freedom
\cite[Section
8.4]{rosner2011biostatistics}.
With unequal variances one uses Welch's
statistic
$$
t = \frac{\bar{x}_1 - \bar{x}_2}
{\sqrt{\frac{s_1^2}{n_1}
+ \frac{s_2^2}{n_2}}}
$$
with the Satterthwaite degrees of freedom
$$
d' = \frac{\left(\frac{s_1^2}{n_1}
+ \frac{s_2^2}{n_2}\right)^2}
{\frac{(s_1^2/n_1)^2}{n_1-1}
+ \frac{(s_2^2/n_2)^2}{n_2-1}}
$$
\cite[Section
8.7]{rosner2011biostatistics}.

\begin{example}
\label{example-two-samples}
In the third exercise sheet the students
generated data in Excel within ranges
given in each problem and compared the
two groups in JMP. With $n = 20$ per
group unless stated:
\begin{itemize}
\item \emph{Latrodectus mactans} (black
widow): time to flee a vibration, in
deciseconds, $11.1$ ($s = 1.62$) at
$1\ \mathrm{Hz}$ versus $2.6$
($s = 1.27$) at $25\ \mathrm{Hz}$;
$t = 18.46$ with $38$ degrees of
freedom, $p < 0.0001$; the variances do
not differ ($F$ test $p = 0.30$).
\item \emph{Crocodylus moreletii}
(Morelet's crocodile): eggs per nest,
$46.1$ ($s = 3.02$) in floating nests in
Veracruz versus $26.1$ ($s = 2.59$) in
nests on aquatic vegetation in San Luis
Potos\'i; $t = 22.45$, $p < 0.0001$.
\item \emph{Galeorhinus galeus} (school
shark): weight $7.81\ \mathrm{kg}$
($n = 21$) in the Gulf of Mexico versus
$9.68\ \mathrm{kg}$ ($n = 19$) in the
Pacific; $t = -4.74$, $p < 0.0001$;
Levene $p = 0.98$.
\end{itemize}
Since the data were simulated from the
stated ranges, these examples train the
method rather than test the biology.
\end{example}

\section{One-way ANOVA and Tukey}
\label{section-anova}

Let $k$ groups have sizes $n_i$, means
$\bar{y}_i$, total size $N$ and grand
mean $\bar{y}$. The total sum of squares
splits as
$$
\sum_{i,j}(y_{ij} - \bar{y})^2
= \sum_i n_i(\bar{y}_i - \bar{y})^2
+ \sum_{i,j}(y_{ij} - \bar{y}_i)^2,
$$
the between and within sums of squares.
Under $H_0\colon \mu_1 = \dots = \mu_k$,
with normal data and equal variances,
$$
F = \frac{\mathrm{SS}_B/(k-1)}
{\mathrm{SS}_W/(N-k)}
\sim F_{k-1,\,N-k}
$$
\cite[Sections
12.2--12.3]{rosner2011biostatistics}.
The ANOVA table lists degrees of freedom,
sums of squares, mean squares
($\mathrm{SS}/\mathrm{df}$), $F$ and $p$;
filling in such tables was a standard
exam exercise.

A significant $F$ says only that some
means differ. Multiple comparison
procedures then find which pairs do while
controlling the overall error rate; the
course used the Tukey-Kramer honestly
significant difference
\cite[Section 12-2]{triola2012elementary}.
Means $i$ and $j$ differ when
$$
|\bar{y}_i - \bar{y}_j| >
q_{1-\alpha;\,k,\,N-k}
\sqrt{\frac{\mathrm{MS}_W}{2}
\left(\frac{1}{n_i}+\frac{1}{n_j}\right)},
$$
with $q$ the studentized range quantile
and $\mathrm{MS}_W = \mathrm{SS}_W/(N-k)$.
JMP prints $q/\sqrt{2}$ as $q^*$. This is
the standard form of the procedure, not
taken from the two textbooks.

\begin{example}
\label{example-anova-lizards}
A lizard exercise compared a control and
two treatments, $20$ animals each, with
means $12.30$, $7.85$ and $10.35$. The
table gave $F(2,57) = 3.88$,
$p = 0.026$, with
$\mathrm{MS}_W = 25.67$; Levene
($p = 0.16$) and Bartlett ($p = 0.35$)
did not reject equal variances. With
$q^* = 2.406$ the Tukey threshold is
$$
2.406\sqrt{25.67\left(\frac{1}{20}
+ \frac{1}{20}\right)} = 3.86,
$$
so only control versus treatment A
(difference $4.45$) is significant.
\end{example}

\section{Simple linear regression}
\label{section-regression}

The model is $y = \beta_0 + \beta_1 x +
\varepsilon$ with independent normal
errors of constant variance. With
$$
\begin{aligned}
L_{xx} &= \sum x_i^2
- \frac{(\sum x_i)^2}{n}, \qquad
L_{yy} = \sum y_i^2
- \frac{(\sum y_i)^2}{n}, \\
L_{xy} &= \sum x_i y_i
- \frac{\sum x_i \sum y_i}{n},
\end{aligned}
$$
the least-squares estimates and the
correlation coefficient are
$$
b_1 = \frac{L_{xy}}{L_{xx}},
\qquad
b_0 = \bar{y} - b_1\bar{x},
\qquad
r = \frac{L_{xy}}{\sqrt{L_{xx}L_{yy}}}.
$$
The regression sum of squares is
$L_{xy}^2/L_{xx}$, the residual sum of
squares is $L_{yy} - L_{xy}^2/L_{xx}$,
and $H_0\colon \beta_1 = 0$ is tested by
$$
F = \frac{L_{xy}^2/L_{xx}}
{(L_{yy} - L_{xy}^2/L_{xx})/(n-2)}
\sim F_{1,\,n-2},
$$
or equivalently by
$$
t = \frac{r\sqrt{n-2}}{\sqrt{1-r^2}}
\sim t_{n-2},
\qquad t^2 = F.
$$
Here $r^2$ is the fraction of the
variance of $y$ explained by the line
\cite[Sections
11.3, 11.4,
11.7, 11.8]{rosner2011biostatistics}.

\begin{example}
\label{example-regression-monkeys}
A worked exercise gave a response
frequency (labelled
$\beta$-$\alpha$ in the sheet) for five
monkeys at ten marijuana doses $x = 0.2,
0.4, \dots, 2.0\ \mathrm{g}$; the mean
responses fell from $31.4$ to $22.4$.
On the ten means, $L_{xx} = 3.3$,
$L_{xy} = -15.84$ and
$L_{yy} = 78.10$, so
$$
\hat{y} = 31.4 - 4.8x,
\qquad r = -0.987,
\qquad r^2 = 0.974,
$$
and $F(1,8) = 294.7$, $p = 1.3 \times
10^{-7}$: the dose explains almost all
the variation in the mean response.
\end{example}

\section{Chi-square test of independence}
\label{section-chi-square}

For an $r \times c$ contingency table
with observed counts $O_{ij}$, row totals
$R_i$, column totals $C_j$ and grand
total $N$, the expected counts under
independence and the test statistic are
$$
E_{ij} = \frac{R_i C_j}{N},
\qquad
X^2 = \sum_{i,j}
\frac{(O_{ij} - E_{ij})^2}{E_{ij}},
$$
and $X^2$ is approximately chi-square
with $(r-1)(c-1)$ degrees of freedom if
the expected counts are not too small
\cite[Section
10.6]{rosner2011biostatistics}.

\begin{example}
\label{example-chi-square-smoking}
Daniel's workbook tests whether
hypertension is independent of smoking:
\begin{center}
\begin{tabular}{lcccc}
 & none & moderate & heavy & total \\
\hline
hypertensive & 21 & 36 & 30 & 87 \\
not & 48 & 26 & 19 & 93 \\
\hline
total & 69 & 62 & 49 & 180 \\
\end{tabular}
\end{center}
The expected count in the first cell is
$87 \cdot 69/180 = 33.35$. Summing over
the six cells gives $X^2 = 14.46$ with
$2$ degrees of freedom, above the
critical value $5.99$ for
$\alpha = 0.05$ ($p = 0.0007$), so
independence is rejected.
\end{example}

\section{Sampling designs}
\label{section-sampling}

The course slides distinguish the
population (elements located in time and
space), the sampling unit, the study
population (units meeting the selection
criteria), the sampling frame, and the
sample. A good sample is representative
and of adequate size; the size depends
on the variability, the precision wanted,
the budget and the population size. An
estimator is a function of the sample
used to estimate a parameter, and its
distribution is the sampling distribution.
The four basic designs
\cite[Section 1-5]{triola2012elementary}:
\begin{itemize}
\item \emph{Simple random}: every subset
of size $n$ has the same chance of being
the sample; there are
$\binom{N}{n}$ possible samples.
\item \emph{Systematic}: a random start,
then every $k$-th unit, with $k = N/n$.
Easy in the field and well spread, but
biased if the population is periodic.
\item \emph{Stratified}: divide the
population into strata homogeneous in the
variable studied and sample within each,
with allocation proportional to stratum
size or optimal (more units from more
variable strata). It increases precision
for a given $n$.
\item \emph{Cluster}: randomly choose
some clusters (groups of units) and
measure all units in them. Cheap when no
frame of individual units exists; best
when clusters are heterogeneous inside
and alike among themselves.
\end{itemize}

In an exam on counts of adult and
juvenile salamanders
(\emph{Rhyacosiredon zempoalensis}) in
several zones, Daniel called the study
observational and descriptive, judged
that counts without zone labels could not
describe the species' distribution, and
proposed a stratified design with zones
of $250\ \mathrm{m}^2$ in which adults
and juveniles are counted separately.

\bibliography{my}
\bibliographystyle{amsalpha}

\end{document}
